\SourceLeaf{85}
\fancyhead[L]{\small The Earth's magnetic force in absolute units}
\renewcommand{\Article}[1]{\par\medskip\Needspace{4\baselineskip}\pdfbookmark[1]{Art. #1}{magnetism-art-#1}\centerline{\textbf{#1.}}\nobreak\smallskip}
\pdfbookmark[0]{The intensity of the Earth's magnetic force expressed in absolute units}{magnetism-body}
\par\medskip
\centerline{THE INTENSITY}
\centerline{OF THE EARTH'S MAGNETIC FORCE}
\centerline{EXPRESSED IN ABSOLUTE UNITS.}
\medskip
The complete determination of the Earth's magnetic force at a given place requires three elements. The first is declination, the angle between the plane in which the force acts and the meridian plane. The second is the inclination of its direction to the horizontal plane. The third is intensity. Declination is to be regarded as the principal element in all applications to navigation and geodesy. It occupied observers and physicists from the very beginning, and for a century they have also devoted constant attention to inclination. The third element, intensity, remained wholly neglected until more recent times, although it is equally worthy of scientific study. Among his many other distinctions, the illustrious \textsc{Humboldt} deserves credit for being among the first to turn his attention to this subject. On his travels he collected many observations of the relative intensity of terrestrial magnetism, which revealed that the intensity increases steadily as we go from the magnetic equator toward the pole. Very many physicists, following in the footsteps of this investigator of nature, have now brought together so many determinations that the distinguished \textsc{Hansteen}, who has done so much for the study of terrestrial magnetism, has recently been able to publish a specimen of a world map showing isodynamic lines.

Their method is to observe how long the same magnetic needle takes to make the same number of oscillations at different places, or to observe how many oscillations it makes in the same interval of time. The intensity is taken to be proportional to the square of the number of oscillations in a given time. In this way we compare total intensities when a dip needle suspended at its center of gravity oscillates about a horizontal axis perpendicular to the magnetic meridian, or intensities of the horizontal force when a horizontal needle oscillates about a \SourceLeaf{86}vertical axis. The second mode of observation is more precise, and its results can readily be converted to total intensities once the inclinations are known.

The method clearly depends on the assumption that the distribution of free magnetism in the particles of the needle used in such a comparison stays the same in each experiment. If the magnetic force of the needle had weakened somewhat over time, the needle would consequently oscillate more slowly afterward. An observer unaware of the change would assign too low a value to the intensity of terrestrial magnetism at the later location. If the experiments cover a moderate interval and the needle is made of well-hardened steel and carefully magnetized, there is little reason to fear an appreciable loss of strength. Moreover, using several needles for the comparisons will reduce the uncertainty. Finally, the assumption gains credibility if, after the journey, the time of oscillation of the needle is found to be unchanged at the first location. But whatever precautions are taken, a gradual weakening of the needle can hardly be avoided, and such agreement can seldom be expected after a longer absence. In comparisons of intensity at widely separated places on Earth, we will therefore usually fall short of the precision and certainty we must seek.

This drawback is less serious when comparing intensities measured at the same time or at times close together. But experience has taught us that both declination and inclination change continuously at a given place, and that these changes become very large after many years. There can be no doubt that the intensity of terrestrial magnetism is subject to similar, almost secular changes. Clearly the method is entirely ineffective for \textit{this} question. Yet for the advancement of natural science it would be highly desirable to shed as much light as possible on this very important question. That certainly cannot be done unless the purely comparative method is replaced by one entirely independent of accidental differences between needles, which expresses the intensity of terrestrial magnetism in fixed units and absolute measures.

It is easy to establish the theoretical principles on which such a long-desired method must rest. The number of oscillations a needle makes in a given time depends both on the intensity of terrestrial magnetism and on the constitution of the needle, namely on the static moment of its elements of free magnetism and on its moment of inertia. Since this moment of inertia can be determined \SourceLeaf{87}without difficulty, observing the oscillations evidently gives us the product of the intensity of terrestrial magnetism and the static moment of the needle's magnetism. The two quantities can be separated only by bringing in observations of another kind that involve a different combination of them. For this purpose introduce a second needle and expose it to both terrestrial magnetism and the magnetism of the first needle, so that the ratio of the two actions can be determined. Both actions will, of course, depend on the distribution of free magnetism in the second needle. The latter action depends in addition on the constitution of the first needle, the distance between the centers, the position of the line joining the centers relative to the magnetic axes of the two needles, and finally on the law of magnetic attractions and repulsions. The immortal \textsc{Tobias Mayer} had already been the first to conjecture that this law agrees with the law of gravitation in that these actions too decrease in inverse proportion to the square of the distance. The experiments of the distinguished \textsc{Coulomb} and \textsc{Hansteen} made this conjecture highly plausible, and the latest experiments put it beyond all doubt. But we must take careful note that this law refers to individual elements of free magnetism. The total effect of a magnetized body is quite different. At very great distances, as can be deduced from that very law, it will approach inverse proportionality to the cube of the distance. Thus, as the distance increases continuously with everything else unchanged, the action of the needle multiplied by the cube of the distance converges asymptotically to a constant. When distances are expressed as numbers using an arbitrary length as unit, that constant has the same dimensions as the action of the Earth's force and can be compared with it. This limiting ratio must be found by suitably arranging and carrying out the experiments. Since it involves only the static moment of the magnetism of the first needle, we then have the quotient obtained by dividing this moment by the intensity of the Earth's force. Comparing it with the product of these quantities found earlier lets us eliminate the static moment and gives the intensity of terrestrial magnetism.

There are two ways to test the actions of terrestrial magnetism and of the first needle on the second needle, since we can observe the second needle either in motion or in equilibrium. In the first method we observe its oscillations while the action of terrestrial magnetism combines with the action of the first needle. The first needle is placed at a suitable distance with its axis in the magnetic meridian through the center of the oscillating needle. The oscillations are thereby accelerated or retarded according to whether friendly or hostile poles \SourceLeaf{88}face one another. We may compare the times of oscillation for the two positions of the first needle with one another, or compare the time for either position with the time of oscillation under terrestrial magnetism alone, with the first needle removed. This gives the ratio of the Earth's force to the action of the first needle. The second method places the first needle so that the direction of the force it exerts in the region of the freely suspended second needle makes an angle, for example a right angle, with the magnetic meridian. The second needle is thereby deflected from the magnetic meridian, and the size of the deflection gives the ratio of the Earth's magnetic force to the action of the first needle.

The first method agrees in essentials with the one the distinguished \textsc{Poisson} proposed some years ago. But the attempts made by some physicists according to his directions, at least those that have come to my attention, either failed altogether or gave only a rough approximation.

The main difficulty arises from the following. From the actions of the needle observed at moderate distances we must calculate a limit that refers to a distance that is, as it were, infinitely great. The more unknowns depending on the individual condition of the needles we have to eliminate, the more the smallest errors of observation disturb, or even ruin, the eliminations needed for this purpose. Yet the number of unknowns can be reduced to very few only if the actions have been observed at distances large enough relative to the length of the needles, and are therefore already very small. The practical means used so far are insufficient to measure such small actions accurately.

I therefore saw that my first task was to devise new means of observing and measuring both the times of oscillation and the directions of the needles with far greater precision than had been possible before. The work undertaken for this purpose and continued for several months, in which the excellent \textsc{Weber} helped me in many ways, achieved the desired result, meeting expectations and going far beyond them. All that is still needed to bring the precision of these experiments up to that of astronomical observations is a place completely sheltered from the influence of nearby iron and from disturbances of the air. Two instruments are available, remarkable as much for their simplicity as for the precision they afford. I must reserve their description for another occasion. In this treatise I present to physicists the experiments made so far in our observatory to determine the intensity of terrestrial magnetism.
